\chapter{補足知識}

\section{\cmd{.venv/} フォルダの扱い}

\cmd{.venv/} は各自の PC で \cmd{uv sync} を実行すれば再生成できます。
以下のことを覚えておいてください。

\begin{itemize}
  \item 共有フォルダに含める必要はない
  \item メールや zip でやり取りする必要はない
  \item 誤って削除しても \cmd{uv sync} で復元できる
  \item リポジトリの \cmd{.gitignore} で除外済み
\end{itemize}

\section{環境を作り直す}

環境が壊れたと感じたときは \cmd{.venv/} を削除して作り直せます。

\begin{wslbox}[環境の再作成]
\begin{lstlisting}[style=bash]
$ cd ~/NMR-lab-kadai/2026_B4_assignment
$ rm -rf .venv
$ uv sync
\end{lstlisting}
\end{wslbox}

\section{新しいパッケージを追加する}

演習の範囲を超えて新しいライブラリを試したい場合：

\begin{wslbox}[パッケージの追加]
\begin{lstlisting}[style=bash]
$ uv add パッケージ名
\end{lstlisting}
\end{wslbox}

例えば \cmd{pandas} を追加する場合：

\begin{wslbox}[pandas の追加例]
\begin{lstlisting}[style=bash]
$ uv add pandas
Resolved 49 packages in 0.3ms
Installed 2 packages in 1.2s
 + pandas==2.2.3
 + python-dateutil==2.9.0
\end{lstlisting}
\end{wslbox}

\cmd{pyproject.toml} と \cmd{uv.lock} が自動的に更新されます。

\section{スクリプトを直接実行する}

ノートブックを使わず Python スクリプト（\cmd{.py} ファイル）を
直接実行したい場合も \cmd{uv run} を使います。

\begin{wslbox}[スクリプトの直接実行]
\begin{lstlisting}[style=bash]
$ cd ~/NMR-lab-kadai/2026_B4_assignment/No.1_FID
$ uv run python makesignal.py
\end{lstlisting}
\end{wslbox}

\section{MATLABとのデータ互換}

このプロジェクトでは \cmd{scipy.io} を通じて
MATLAB の \cmd{.mat} ファイルを読み書きできます。

\begin{lstlisting}[style=python, caption={.mat ファイルの読み込み例}]
import scipy.io

# MATLAB の .mat ファイルを読む
data = scipy.io.loadmat("Signal.mat")

# Python の配列を .mat ファイルに保存する
scipy.io.savemat("result.mat", {"signal": signal})
\end{lstlisting}

\section{インストール済みパッケージの確認}

\begin{wslbox}[パッケージ一覧の表示]
\begin{lstlisting}[style=bash]
$ uv pip list
\end{lstlisting}
\end{wslbox}
